% #############################################################################
% Abstract Text
% !TEX root = ../main.tex
% #############################################################################
% reset acronyms
\acresetall
% use \noindent in firts paragraph
\noindent Machine-learning-based network intrusion detectors can be evaded by reshaping only packet timing and sizes, while preserving protocol compliance and attack functionality.
Yet many defenses are evaluated on feature perturbations that may not be realizable as traffic.
This thesis proposes DEPT, a streaming defense that detects and purifies manipulated traffic before it reaches a protected detector.

DEPT represents traffic as masked packet-level bits and flow statistics at two temporal granularities.
At each packet, an attention-based meta-learner fuses a Random Forest and a Temporal Convolutional Network to detect manipulated flows.
Flagged flows are corrected in feature space by a causal purification Transformer and mapped back to packet timing and sizes.

TANTGen~\cite{pissarra2026tantgen}, a black-box generator, supplies the adversarial traffic for training and evaluation: ten sources from three public datasets, each in three perturbation modes.
Manipulation detection reaches an end-of-flow AUROC of $0.933$; the Random Forest alone scores higher from the second packet, but the ensemble predicts from the first and better resists the evaluated attacks on the sequence model.
On the independent detector each attack targeted, purification recovers $63\%$ of the detection lost to manipulation.
End to end, DEPT cuts attack success from $21.1\%$ to $11.6\%$, whereas the diffusion-based purifier of Merzouk et al.~\cite{merzouk2024diffusion} recovers almost none of the lost detection.
Of DEPT's realizable packets, $97.2\%$ parse without malformed-packet errors.
Results vary by source: manipulation detection is near or below chance on size-perturbed reflection floods, and purifying flows that did not evade can lower downstream detection.
